\documentclass[10pt,a4paper]{amsart}
\usepackage[margin=19mm]{geometry}
\usepackage{amsmath,amssymb,mathtools}
\usepackage[scr=rsfs,cal=euler]{mathalfa}
\usepackage{xfrac}
\usepackage[unicode,hidelinks,bookmarks=false]{hyperref}
\newcommand{\ie}{i.e.\ }
\newcommand{\cf}{cf.\ }
\newcommand{\resp}{respectively\ }
\newcommand{\Subsection}{\subsection}
\providecommand{\nobreakdash}{\nobreak\hbox{-}}
\setlength{\emergencystretch}{2em}
\multlinegap0pt
\usepackage{bm}
\usepackage{bbm}
\usepackage{dsfont}
\usepackage{xspace}
\usepackage{cancel}
\usepackage{mathtools}
\usepackage{amsthm}
\makeatletter
\@ifundefined{lemma}{\newtheorem{lemma}{Lemma}}{}
\@ifundefined{proposition}{\newtheorem{proposition}{Proposition}}{}
\makeatother
\newcommand{\Cbar}{\underline{C}}
\newcommand{\be}{\begin{equation}}
\newcommand{\ee}{\end{equation}}
\newcommand{\hba}{\underline{h}}
\newcommand{\hsp}{\hspace{.5mm}}
\newcommand{\upr}{u^{\prime}}
\title[Formation of trapped surfaces for the spherically symmetric ]{Formation of trapped surfaces for the spherically symmetric Einstein-Yang-Mills system with non-trivial incoming data}
\author{Nikolaos Athanasiou, Puskar Mondal, Shing-Tung Yau}
\date{Source: arXiv:2608.09746v1 (CC BY 4.0)}
\begin{document}
\maketitle

\noindent\emph{Source and licence.} Formation of trapped surfaces for the spherically symmetric Einstein-Yang-Mills system with non-trivial incoming data. Version arXiv:2608.09746v1 (CC BY 4.0), distributed under the Creative Commons Attribution 4.0 International licence, as recorded in the arXiv licence metadata for this exact version and shown on the abstract page https://arxiv.org/abs/2608.09746v1. Full rights notice: licenses/arxiv-2608.09746-CC-BY-4.0.txt.\par\smallskip

\noindent\textit{Editorial scope.} Counted unit: the Lemma whose source label is \texttt{None} in the article (source0.tex lines 867--871, source label \texttt{None}), delivered together with the article's own complete original proof of it (source0.tex lines 873--927, one \texttt{proof} environment of 55 lines). No mathematics is written, repaired or re-derived: statement and proof are the article's own text line for line. Required context retained uncounted: propmonotonicity (lines 751--753); lemmaomegaeta (lines 784--788); Thetalemma (lines 798--802); hlemma (lines 837--841). Every definition, display and label of those lines is retained unchanged. Omitted from this package and not redistributed: 1--750, 754--783, 789--797, 803--836, 842--866, 872--872, 928--1066. Nothing inside a retained excerpt is omitted or altered: every retained body is delivered exactly as the article's lines are written, so no omission receipt of any kind accompanies this package. Citations: the retained excerpts carry no citation command at all, so no bibliography entry of the article is transcribed and no reference is redistributed. Reference adaptation: none was needed and none was made; every reference inside the delivered unit resolves inside it, so no unresolved reference or citation is produced. Printed slips kept literal: no printed word, symbol, hypothesis or number of the counted statement or of its proof was altered. Layout and class: the article's own class is \texttt{article} with packages including graphicx, hyperref, xcolor, amsthm, tikz, amsmath, amssymb; the wrapper is the lane's pinned \texttt{amsart} editorial wrapper at 10pt on A4, which restates the theorem-like environments the retained text needs and adds the article's own macro definitions that the retained text needs (\texttt{\textbackslash Cbar}, \texttt{\textbackslash be}, \texttt{\textbackslash ee}, \texttt{\textbackslash hba}, \texttt{\textbackslash hsp}, \texttt{\textbackslash upr}), with extra wrapper packages (bm, bbm, dsfont, xspace, cancel, mathtools). The delivered numbering is the unit's own. Assets: no retained span contains a figure environment, a graphics-inclusion command, a PDF-page inclusion, a table or a TikZ picture; no image file is copied into this package, the package declares no asset dependency and no image file is redistributed with it. Licence: version arXiv:2608.09746v1 of this article is distributed under the Creative Commons Attribution 4.0 International licence, as recorded in the arXiv licence metadata for this exact version and shown on the abstract page https://arxiv.org/abs/2608.09746v1. A full rights notice is delivered at licenses/arxiv-2608.09746-CC-BY-4.0.txt. Characters: every retained excerpt is pure ASCII, so the delivered engine, XeLaTeX with its default Latin Modern text font, reports no missing character.
\begin{proposition} \label{propmonotonicity}
Assume that the initial data along $\Cbar$ is magnetically subextremal and that $\mathcal{R}$ is free of trapped surfaces. Then $\partial_u\partial_v r \leq 0$ in $[u_0, u_*] \times [v_1,v_2]$ and $\delta(u):= \frac{r_2(u)}{r_1(u)}-1 \leq \frac{1}{2}$ for $u \in [u_0, u_*],$ where $u_*$ is defined so that $x(u_*)= \frac{3\delta_0}{1+\delta_0}.$ 
\end{proposition}

\begin{lemma} \label{lemmaomegaeta}
Assume that the initial data along $\Cbar$ is magnetically subextremal and that $\mathcal{R}$ is free of trapped surfaces. Then for all $(u,v_a) \in [u_0, u_*] \times [v_1,v_2]$, we have the following estimate for the charge:

\[  \frac{Q_a^2}{r_a^2}\leq \frac{\omega}{4}\eta_a+2\varepsilon,  \]where we recall that $\varepsilon = \sup_{C \cup \Cbar}\frac{Q^2}{r^2} <1$. If, moreover, $\eta_a \geq \frac{8\varepsilon}{\omega},$ then $\frac{Q_a^2}{r_a^2}\leq \frac{\omega}{2}\eta_a$.
\end{lemma}

\begin{lemma} \label{Thetalemma}
Define $\Theta:= \partial_u w_2 - \partial_u w_1$. Suppose that the initial data along $\Cbar$ is not super-charged and that $\mathcal{D}(0,v_1)\cup \mathcal{R}$ is free of trapped surfaces. If $\eta \geq \frac{8 \varepsilon}{\omega}$, then 

\[  \Theta(u)^2 \leq \big(1 +\frac{\omega}{2}\big) \frac{- \partial_u r_2}{\partial_v r_2}(m_2-m_1)\bigg( \frac{1}{r_1}- \frac{1}{r_2}\bigg)(u),    \]for all $u \in [u_0, u_*]$.
\end{lemma}

\begin{lemma} \label{hlemma}
Assume that $\eta \geq \frac{8 \varepsilon}{\omega}$, the initial data along $\Cbar$ is not magnetic-supercharged and that $\mathcal{D}(0,v_1)\cap \mathcal{R}$ is devoid of trapped surfaces. Then

\[ \frac{h_2(u)}{h_1(u)}\leq e^{-(1-\frac{\omega}{2})}\eta(u).      \]
\end{lemma}

\begin{lemma}
Assume that the region $\mathcal{D}(0,v_1)\cap \mathcal{R}$ is free of trapped surfaces and the initial data along $\Cbar$ is magnetically subextremal. Then, if $\eta_0 \geq \frac{13\varepsilon}{\omega}+g_{\omega}(\delta_0)$, we have $\eta(x) \geq \frac{12\varepsilon}{\omega}$ for all $x= x(u)\in [\frac{3\delta_0}{1+\delta_0},1]$. Here, $g_{\omega}(x)$ is defined as:

\[ g_{\omega}(x):= \frac{1+\frac{\omega}{2}}{1-\frac{\omega}{2}} \hsp \frac{1}{(1+x)^2}\bigg(\bigg( \frac{2^{1-\frac{\omega}{2}}}{\omega}+ \frac{1}{2^{1+\frac{\omega}{2}}(1+\frac{\omega}{2})}\bigg)x^{1-\frac{\omega}{2}}-\frac{2}{\omega}x-\frac{1}{1+\frac{\omega}{2}}x^2\bigg).    \]
\end{lemma}

\begin{proof}
Let $\upr:= \sup\{ u \in [u_0,u_*] \mid \eta(s) \geq \frac{12\varepsilon}{\omega} \hsp \text{for} \hsp \text{all} \hsp s\in [u_0,u]\}.$ The quantity we will estimate is $\frac{\text{d}\eta}{\text{d}x}$, where $x= \frac{r_2(u)}{r_2(u_0)}$. There holds 

\begin{align} \label{coreinternal} \frac{\text{d}\eta}{\text{d}x} =& \frac{\frac{\text{d}\eta}{\text{d}u}}{\frac{\text{d}x}{\text{d}x}}= \frac{r_2(u_0)}{\hba_2}\big(-\frac{2\hsp \hba_2}{r_2^2}(m_2-m_1) + \frac{2}{r_2}\partial_u (m_2-m_1)\big) \notag \\ =& -\frac{\eta}{x}+\frac{2}{x \hsp \hba_2} \bigg(\frac{\hba_2 \hsp Q_2^2}{2r_2^2}-\hba_2\hsp (\partial_u w_2)^2 - \frac{\hba_1 \hsp Q_1^2}{2\hsp r_1^2}+h_1(\partial_u w_1)^2\bigg) \notag \\ \leq& -\frac{\eta}{x}-\frac{2\hsp h_2}{x \hsp \hba_2}\hsp \big((\partial_u w_2)^2 - \frac{h_1}{h_2}(\partial_u w_1)^2\big) + \frac{Q_2^2}{x\hsp r_2^2}. \end{align} 
We now focus our attention on the region $[u_0, \upr]$. Since $\eta \geq \frac{12\varepsilon}{\omega}$ by the bootstrap asumption in $[x^\prime,1]$, there also holds $\eta \geq \frac{8\varepsilon}{\omega}$, whence the conditions of Lemma \ref{hlemma} hold. We thus obtain:

\begin{align}
(\partial_u w_2)^2 - \frac{h_1}{h_2}(\partial_u w_1)^2 \leq& (\partial_u w_2)^2 - e^{\eta(1-\frac{\omega}{2})} (\partial_u w_1)^2\notag \\ =& \Theta^2 +2\Theta \hsp (\partial_u w_1) + \big(1-e^{\eta(1-\frac{\omega}{2})}\big) (\partial_u w_1)^2. \label{intermediate1}
\end{align}The expression above on the right-hand side of \eqref{intermediate1} is then bounded as follows:

\begin{align}
\Theta^2 +2\Theta \hsp (\partial_u w_1) + (1-e^{\eta(1-\frac{\omega}{2})}) (\partial_u w_1)^2 \leq& \big(1 + \frac{1}{e^{\eta(1-\frac{\omega}{2}}-1}\big) \Theta^2 \notag \\ \leq& \big(1 + \frac{1}{\eta(1-\frac{\omega}{2})}\big) \Theta^2,
\end{align}since $\eta(1-\frac{\omega}{2}) \geq 0$. Plugging this bound back to \eqref{coreinternal}, we have

\begin{align}
\frac{\text{d}\eta}{\text{d}x}\leq -\frac{\eta}{x}- \frac{2h_2}{x \underline{h}_2}\big(1 + \frac{1}{\eta(1-\frac{\omega}{2})}\big) \Theta^2 + \frac{Q^2}{xr_2^2}. 
\end{align}

\noindent Applying Lemma \ref{Thetalemma}, we have

\begin{equation}\label{mediumestimate} \frac{\text{d}\eta}{\text{d}x}\leq -\frac{\eta}{x} + \frac{\eta}{x}\big(\frac{\omega}{2}+1\big)\bigg(1+ \frac{1}{\eta(1-\frac{\omega}{2})}\bigg) \big(\frac{r_2}{r_1}-1\big) + \frac{Q_2^2}{xr_2^2}. \end{equation}By employing the monotonicity formula from Proposition 
\ref{propmonotonicity}, we get

\begin{align}\label{boundmonot}
\delta(u)= \frac{r_2(u)-r_1(u)}{r_2(u)-(r_2(u)-r_1(u))}\leq& \frac{r_2(u_0)-r_1(u_0)}{r_2(u_0)-(r_2(u_0)-r_1(u_0))} \notag \\ =& \frac{\delta_0}{x(u)(1+\delta_0)-\delta_0}. 
\end{align}The last term on the right-hand side of \eqref{mediumestimate} will be bounded through the use of Lemma \ref{lemmaomegaeta}. Looking at \eqref{mediumestimate} and using \eqref{boundmonot} together with Lemma \ref{lemmaomegaeta} (which is applicable since $\eta \geq \frac{12\varepsilon}{\omega}>\frac{8\varepsilon}{\omega}$ in our region of interest),

\begin{align}
\frac{\text{d}\eta}{\text{d}x}\leq& \eta \bigg(\big(1+\frac{\omega}{2}\big) \frac{\delta_0}{x^2(1+\delta_0)-x\delta_0)}-\frac{1}{x}\bigg) +\frac{1+\frac{\omega}{2}}{1-\frac{\omega}{2}} \frac{1}{x}\frac{\delta_0}{x(1+\delta_0)-\delta_0}+\frac{Q_2^2}{xr_2^2} \notag \\ \leq&\eta \bigg(\big(1+\frac{\omega}{2}\big) \frac{\delta_0}{x^2(1+\delta_0)-x\delta_0)}-\frac{1}{x}\bigg) +\frac{1+\frac{\omega}{2}}{1-\frac{\omega}{2}} \frac{1}{x}\frac{\delta_0}{x(1+\delta_0)-\delta_0}+\frac{\omega}{2}\frac{\eta}{x} \notag \\ =& -\frac{\eta}{x} \bigg(1-\frac{\omega}{2} - (1+\frac{\omega}{2}) \frac{\delta_0}{x(1+\delta_0)-\delta_0} \bigg) +\frac{1+\frac{\omega}{2}}{1-\frac{\omega}{2}} \frac{1}{x}\frac{\delta_0}{x(1+\delta_0)-\delta_0}.\notag
\end{align}It thus follows that if we define

\[ g(x):=   1-\frac{\omega}{2} - (1+\frac{\omega}{2}) \frac{\delta_0}{x(1+\delta_0)-\delta_0}  \]and \[f(x) :=\frac{1+\frac{\omega}{2}}{1-\frac{\omega}{2}} \frac{1}{x}\frac{\delta_0}{x(1+\delta_0)-\delta_0},  \]we obtain the following differential inequality for all $x\in [x^{\prime},1]$:

\[ \frac{\text{d}\eta}{\text{d}x}+\eta \frac{g(x)}{x}-\frac{f(x)}{x}\leq 0.    \]A straightforward calculation gives, solving the ODE, the bound

\be \eta_0 - e^{-G(x)}\eta(x)\leq F(x), \label{almostlast}\ee

where \be G(x) := \int_x^1 \frac{g(t)}{t} \hsp \text{d}t= \ln \bigg(\frac{(x(1+\delta_0)-\delta_0)^{1+\frac{\omega}{2}}}{x^2}\bigg)\ee and

\begin{align}
&F(x) := \int_x^1 e^{-G(t)}\frac{f(t)}{t} \text{d}t \notag \\ =& \frac{1+\frac{\omega}{2}}{1-\frac{\omega}{2}} \hsp \frac{\delta_0}{(1+\delta_0)^2}\bigg(\frac{2}{\omega}\hsp \frac{1}{(x(1+\delta_0)-\delta_0)^{\frac{\omega}{2}}}-\frac{2}{\omega}+\frac{1}{1+\frac{\omega}{2}}\hsp\frac{\delta_0}{(x(1+\delta_0)-\delta_0)^{\frac{\omega}{2}}}- \frac{\delta_0}{1+\frac{\omega}{2}}\bigg). 
\end{align}Since $F$ is monotonically decreasing, it attains its maximum in the interval $\big[ \frac{3\delta_0}{1+\delta_0},1\big]$ at $x= \frac{3 \delta_0}{1+\delta_0}$. Therefore, 

\begin{align}
F(x)\leq F\big(\frac{3\delta_0}{1+\delta_0}\big)= g_{\omega}(\delta_0).
\end{align}Substituting into \eqref{almostlast}, we get

\begin{align} \label{last}\notag \eta(x)\geq e^{G(x)}(\eta_0-F(x))=& \frac{(x(1+\delta_0)-\delta_0)^{1+\frac{\omega}{2}}}{x^2}(\eta_0-F(x)) \\ \geq& \frac{(x(1+\delta_0)-\delta_0)^{1+\frac{\omega}{2}}}{x^2}(\eta_0-g_{\omega}(\delta_0)). \end{align}Since $\omega<\frac{2}{3}$, it is a straightforward calculation that 

\[  \sup_{x \in[x_*,1]} \frac{x^2}{\big(x(1+\delta_0)-\delta_0\big)^{1+\frac{\omega}{2}}}=1.   \]Taking into account the hypothesis 

\[ \eta_0 > \frac{13 \varepsilon}{\omega}+g_{\omega}(\delta_0),   \]we obtain the inequality

\[  \eta_0 > \frac{13\varepsilon}{\omega}+g_{\omega}(\delta_0) \geq \frac{13 \varepsilon}{\omega}  \frac{x^2}{\big(x(1+\delta_0)-\delta_0\big)^{1+\frac{\omega}{2}}} +g_{\omega}(\delta_0), \]for $x\in[x^{\prime},1]$. Substituting the above into \eqref{last}, we obtain \[ \eta(x) \geq \frac{13\varepsilon}{\omega}, \hspace{2mm} \text{for all} \hspace{1mm} x\in[x^{\prime},1].    \]However, by continuity of $\eta,$ we can find $\tilde{x}<x^{\prime}$ such that for all $x \in[\tilde{x},1]$ there holds $\eta(x)\geq \frac{12\varepsilon}{\omega}$, or equivalently $\eta(u)\geq \frac{12\varepsilon}{\omega}$ for all $u \in [u_0,\tilde{u}]$, contradicting the supremum property of $u^{\prime}$.
\end{proof}

\end{document}
